\section{ADR-011: DAW session export}
\textbf{Status:} accepted and implemented (future work, FR-DAW-03).

\textbf{Context.} A cover musician who records over the stems wants them in a
DAW project, not as loose files: every stem on its own track starting at 0, the
mixer's balance, and a bar grid that matches the song so the metronome, loops
and quantising work. M14 already stores the tempo, beat grid, section markers
and key of every song (ADR-010).

\textbf{Decision.} \texttt{POST /items/\{id\}/session} writes a new folder
\texttt{<title> - DAW session} (never overwriting) with:
\begin{itemize}
  \item the audio: every stem, the \emph{Residual} (the mix minus every stem;
  muted when the mixer drops it) and the \emph{Original mix} (muted, for A/B),
  as 32-bit float WAV named like the stem export
  (\texttt{<title> - Vocals.wav}). WAV keeps every file sample-aligned at 0;
  MP3 would add encoder delay;
  \item \texttt{<title>.RPP}: REAPER's documented plain-text format, written
  line by line (track \texttt{VOLPAN} and \texttt{MUTESOLO}, one item per track
  at position 0, \texttt{MARKER} lines, a \texttt{TEMPOENVEX} chunk with square
  tempo points and time signatures);
  \item \texttt{<title>.dawproject}: the open DAWproject format (version 1.0,
  specification and XML Schema under the MIT licence by Bitwig), read by
  Bitwig Studio, Studio One and Cubase. The arrangement uses a seconds
  timeline, so clips, markers and tempo points sit exactly where the audio has
  them; the audio is referenced next to the file (\texttt{external="true"})
  instead of embedded, so the formats share one copy;
  \item \texttt{<title> - tempo map.mid}: a Standard MIDI File (format 0, meta
  events only: tempo, time signature, key signature, markers) for any other
  DAW: import it and drop the WAV files at 0.
\end{itemize}
All three writers are pure Python with no third-party code. Fader gain, pan and
mute (after solo) come from the mixer; pitch and tempo changes from the
practice tools are not applied, because the session is for recording over the
original.

\textbf{Tempo map.} The audio is never moved or stretched, so the grid is fitted
to it. If a straight bar grid is within 40~ms of every detected downbeat (a
click-tracked recording; the beat tracker works on 23~ms frames) the song gets
one tempo; otherwise one tempo per bar, so every bar line falls on a downbeat.
The time before the first downbeat becomes a pickup bar of whole beats in its
own time signature (for example 1/4), so the following bar lines land on the
downbeats; a pickup shorter than one beat at 960~BPM (REAPER's limit) joins the
first bar. Neither RPP nor DAWproject has a key signature field, so the key goes
in the REAPER project notes and the DAWproject metadata comment; the MIDI file
carries it as a key signature event.

\textbf{Alternatives.} Embedding the audio in the \texttt{.dawproject} ZIP
makes one portable file but stores every stem twice next to the RPP. A beats
timeline for DAWproject (what Bitwig writes) would need the audio warped to the
tempo map. Ableton Live, Logic Pro and Pro Tools projects are proprietary and
undocumented, so they are reached through the MIDI tempo map. No library was
needed: the only DAWproject library is Java, and the RPP Python parsers read
rather than write.

\textbf{Consequences.} Tests compare each writer with golden files built from a
synthetic session and validate the DAWproject XML against the official
\texttt{Project.xsd} and \texttt{MetaData.xsd} (copied into the tests, MIT,
checked with \texttt{xmlschema}, a development-only dependency). The projects
reference their audio by relative path, so the folder must stay together.
Pan is written as a balance value; a DAW's own pan law may make an off-centre
stem slightly louder or softer than in the preview. Songs in other metres get
4/4 bars, as in ADR-010.
